%% SimulacraBench method description -- draft 7 (2026-11-19), in the organisers'
%% report skeleton (simulacrabench-template.sty from the OpenReview form, kept as
%% organisers-template.tex.txt) with the official neurips_2026.sty [final].
%% Build: bash report/build.sh  (writes _report_build/method.pdf; the author line
%% comes from a file the build script writes outside Git).

\documentclass{article}
% The organisers' template loads [final] only; neurips_2026.sty (2026-01-29) then leaves
% \@trackname undefined, so the default track option `main' is named explicitly.
\usepackage[main,final]{neurips_2026}

\usepackage[utf8]{inputenc}
\usepackage[T1]{fontenc}
\usepackage{hyperref}
\usepackage{url}
\usepackage{booktabs}
\usepackage{amsfonts}
\usepackage{amsmath}
\usepackage{amssymb}
\usepackage{nicefrac}
\usepackage{microtype}
\usepackage{xcolor}
\usepackage{graphicx}
\usepackage{enumitem}
\usepackage{seqsplit}
\usepackage{float}
\usepackage[htt]{hyphenat}
\setlength{\emergencystretch}{3em}
% Filled from the Test-phase probe (submission 1054224, failed on purpose 2026-11-16, status_details read from our own submission API).
\newcommand{\ProbeType}{NotImplementedError}
\newcommand{\ProbeMeaning}{on the grader's GPU TabICL found its checkpoint, completed its fold passes and its full-data pass for at least one item, and was mixed into at least one item, and all three instruments were computed within 75\% of the 3600\,s budget (the run ended about 25 minutes after upload). The failed probe did not use the single Test submission}
\newcommand{\TabICLUsed}{TabICLv2, \texttt{jingang/TabICL} (BSD-3-Clause), checkpoint \texttt{tabicl-classifier-v2-20260212.ckpt}, revision \texttt{\seqsplit{4dcd344ece2c00be9e831fdd35bed57b5ad83e19}}}
\newcommand{\CodeURL}{https://github.com/junhyun-dev/simulacrabench}
   % \ProbeType, \ProbeMeaning, \TabICLUsed, \CodeURL

\newcommand{\teamid}{\texttt{junhyun-dev}}

\newcommand{\cbox}{$\square$}
\newcommand{\cboxchecked}{$\boxtimes$}

\title{SimulacraBench Technical Report:\\ Skip-aware backoff with out-of-fold mixing}

\InputIfFileExists{_author.tex}{}{\author{Team junhyun-dev}}

\begin{document}

\maketitle

\begin{center}
\begin{tabular}{l}
\toprule
\textbf{Team ID:} \teamid \\
\bottomrule
\end{tabular}
\end{center}

\begin{abstract}
We complete each held-out respondent's survey from the only answers the grader
shows for them, the GIVEN block. For every PREDICT item we fit five predictors
on the visible rows inside the scoring container: a hierarchical Dirichlet
backoff over the most informative GIVEN columns, a gate recursion that carries a
parent's predicted answer down a skip chain, a multitask network, per-item
gradient-boosted trees, and the pretrained TabICLv2 tabular model used in
context without fine-tuning. Per item we choose between and mix them by
out-of-fold log score computed the way the grader computes it, then choose a
temperature. Hidden rows are predicted by models refitted on every visible row.
We use no language model, no question text and no external data. Our Test-phase
entry scores $0.48$ skill ($0$ = crowd marginal, $1$ = perfect; $0.34$ UNICEF,
$0.23$ World Bank, $0.87$ UNHCR). Because the grader returns no logs, one
failing probe submission, whose exception type reports only the state of our own
code, confirmed that TabICLv2 ran on the grader's GPU.
\end{abstract}

%% ---------------------------------------------------------------------------
\section{Method summary}
\label{sec:method}

\subsection{Method type}
Marked categories; a one-line description of each marked component follows in Section~\ref{sec:pipeline}.

\begin{table}[h]
  \centering
  \footnotesize
  \begin{tabular}{@{}c@{\hspace{3pt}}p{0.44\linewidth}c@{\hspace{3pt}}p{0.44\linewidth}@{}}
    \toprule
    \cbox & LLM-simulated respondents (prompting, personas) & \cbox & Fine-tuned language model \\
    \cbox & Tabular / conditional generative model (e.g.\ VAE, diffusion, GAN, CTGAN) & \cboxchecked & Supervised prediction (GBDT, neural net, linear/GLM) \\
    \cboxchecked & Bayesian / hierarchical / multilevel model & \cbox & Small-area estimation or post-stratification (e.g.\ MRP) \\
    \cbox & Reweighting, raking, or calibration of the released sample & \cbox & Imputation / matching / nearest-neighbour \\
    \cbox & Retrieval over released or external microdata & \cbox & Agent-based or mechanistic simulation \\
    \cboxchecked & Ensemble or stacking of multiple components & \cboxchecked & Other: pretrained in-context tabular classifier \\
    \bottomrule
  \end{tabular}
\end{table}

\subsection{Pipeline}
\label{sec:pipeline}
Each blank cell is scored by the log of the probability given to the answer that
was there, after the grader renormalises the vector and mixes it with a flat
vector at floor $\varepsilon=10^{-3}$; skill is relative to the crowd marginal.
For each instrument, \texttt{predict()} takes the frame the grader passes in,
splits the visible rows (all PREDICT answers present) into three folds, and for
every PREDICT item: (i)~fits the \emph{hierarchical backoff} (Bayesian
component: a Dirichlet-smoothed chain of conditional distributions over GIVEN
columns) and, for items gated on a PREDICT parent, the \emph{gate recursion};
(ii)~fits a \emph{multitask network} and per-item \emph{gradient-boosted trees}
(supervised prediction); (iii)~runs \emph{TabICLv2}~\citep{qu2025tabicl,qu2026tabiclv2}
in context (other: pretrained in-context tabular classifier); (iv)~mixes them
greedily by out-of-fold log score (ensemble) and chooses a temperature; and
(v)~predicts the hidden rows with the same choices, from models refitted on every
visible row.

\subsection{Inputs and feature construction}
The method uses only the GIVEN columns of each respondent (7, 9 and 17 columns
for UNICEF, World Bank and UNHCR) and, from the schema, each item's option list
and declared gate parent. Hidden respondents show only their GIVEN block; every
PREDICT cell of a hidden row is blank, so other answers are never available at
prediction time. All GIVEN columns are treated as categorical (one-hot for the
network). Question text and other metadata are ignored: on the World Values
Survey, ``Random Forests demonstrate stronger performance than LLMs when trained
with sufficient data''~\citep{sinacola2025llms}, and on SHARE
``demographics-only agents (1) exhibited a central tendency bias, skewing
answers toward population means''~\citep{garzon2026demographics}; with thousands
of visible respondents per instrument we spent the budget on tabular models.
Survey weights are not used. When an item's gate parent is itself PREDICT (155
of the 156 gated UNHCR items), whether it was asked is uncertain and chains up
to four levels deep, so the probability of \texttt{NA\_GATED} is predicted, not
copied.

\subsection{Model or simulator}
\paragraph{Backoff.} For item $j$, with an option list that includes
\texttt{NA\_GATED} when the item is gated, GIVEN columns are ranked by mutual
information with the answer on the training fold. Level $0$ is the marginal
with half a count per option; level $l$ conditions on the first $l$ columns,
$p_l(\cdot\mid s_l)=\bigl(n(s_l,\cdot)+\alpha\,p_{l-1}(\cdot\mid s_{l-1})\bigr)/\bigl(N(s_l)+\alpha\bigr)$,
so a stratum never seen keeps its coarser estimate. When the gate parent is
GIVEN, that column is forced first, so a skip the visible answer decides is
reproduced up to the $\alpha$ smoothing.

\paragraph{Gate recursion.} For an item gated on a PREDICT parent,
$P(\text{asked}\mid x)=\sum_v P_{\text{parent}}(v\mid x)\,g(v)$, where $g(v)$ is
the share of visible respondents with parent answer $v$ who were asked, smoothed
by half a pseudo-count toward the declared rule; the answer given asked is a
backoff fitted on the asked rows. Parents are processed first, so their
out-of-fold and hidden-row predictions feed their children.

\paragraph{Network, trees, TabICLv2.} One network per instrument: one-hot GIVEN
input, a 256-unit hidden layer with dropout plus a zero-initialised linear path
from input to logits (so it contains multinomial logistic regression), and one
softmax per item. Per item, scikit-learn~\citep{pedregosa2011sklearn}
\texttt{HistGradientBoostingClassifier} with categorical GIVEN features.
TabICLv2 (\texttt{jingang/TabICL}, checkpoint
\texttt{tabicl-classifier-v2-20260212.ckpt}) with four estimators and no
fine-tuning: for each fold the other folds' rows (at most 20{,}000) are the
context. For an item whose gate parent is GIVEN, the asked/not-asked split is
imposed on the network, tree and TabICL outputs from the smoothed share of
visible respondents asked for each parent value. Adding TabICLv2 followed the
public README of \texttt{exochard/simulacrabench-v33} (Apache-2.0, read
2026-11-13), which describes a similar per-item out-of-fold blend with
TabICLv2; we read only its README and took no code from it.

\subsection{Fitting, training, or prompting}
Everything is fitted inside \texttt{predict()} on the visible rows of the frame
it receives. The README asks that training happen offline, but the respondent
data are never released (FAQ~4), and in the Test phase the DEV respondents
return as visible rows, ``answers included, which is where they are worth
most''; the official crowd-marginal baseline also counts inside
\texttt{predict()}. No external data. The backoff's $\alpha$ and depth are
chosen by out-of-fold log score; the network uses cross-entropy over all
observed cells, AdamW and early stopping on the held-out fold (three seeds per
fold); the trees stop early on an internal 15\% split. Settings are in
Appendix~\ref{app:config}.

\subsection{Prediction and aggregation}
Per item: pick the backoff or recursion setting with the best floored
out-of-fold log score; then mix in the network with weight $w\in\{0.1,\dots,1.0\}$
if some $w$ improves that score, then the trees, then TabICL, each greedily on
top of the previous mixture. Every comparison uses the grader's own transform
(renormalise, floor, log), so a component that does not help an item is not
mixed into it. Hidden rows use the same choices; the backoff is refitted on all
visible rows, the network is averaged over its fold models, and the trees and
TabICL are refitted once more on every visible row (TabICL keeps the fold
average if time runs out before that pass).

\subsection{Post-processing and calibration (if applicable)}
A per-item temperature $t\in\{0.7,0.8,0.9,1.1,1.25,1.5\}$ for $p^t$
(renormalised) is applied if it improves the out-of-fold score on the visible
rows; no other post-processing.

%% ---------------------------------------------------------------------------
\section{Models, data, and compute}
\label{sec:resources}

\paragraph{Pretrained models and off-the-shelf components.} Table~\ref{tab:models}.

\begin{table}[h]
  \centering
  \small
  \caption{Pretrained models and off-the-shelf components.}
  \label{tab:models}
  \begin{tabular}{@{}p{0.2\linewidth}p{0.48\linewidth}p{0.26\linewidth}@{}}
    \toprule
    Component & Version / checkpoint & Role in pipeline  \\
    \midrule
    TabICLv2 (BSD-3-Clause) & \texttt{jingang/TabICL}, \texttt{\seqsplit{tabicl-classifier-v2-20260212.ckpt}}, revision \texttt{4dcd344e}; \texttt{tabicl==2.2.0} & in-context classifier, one of five mixed predictors \\
    scikit-learn & \texttt{HistGradientBoostingClassifier}, as preinstalled in the grader image & per-item trees \\
    \bottomrule
  \end{tabular}
\end{table}

\paragraph{Data.} Development used the Development phase through Codabench
submissions and the organisers' synthetic practice data
(\texttt{make\_sandbox.py}) to check that the code runs. No external data
(Table~\ref{tab:external}).

\begin{table}[h]
  \centering
  \small
  \caption{External datasets used.}
  \label{tab:external}
  \begin{tabular}{llll}
    \toprule
    Dataset & Source / URL & Size & Use in pipeline \\
    \midrule
    None & -- & -- & -- \\
    \bottomrule
  \end{tabular}
\end{table}

\paragraph{Compute.} The grader's single GPU within the 3600\,s Test budget
(900\,s in Development). Wall-clock guards stop \emph{starting} expensive
components and fall back to what is already fitted: trees after 2016\,s, TabICL
after 2160\,s; after 2520\,s the network stops adding seeds and any item not yet
reached gets the crowd marginal. TabICL runs only when a GPU is visible and the
checkpoint is found. Because the grader returns only an exception's type, we
submitted once a probe identical to the final code that raised an exception
whose type, one of eleven fixed values, reports only the state of our own code;
it returned \texttt{\ProbeType}: TabICLv2 found its checkpoint, ran on the GPU
and was mixed into at least one item, within 75\% of the budget. The competition
terms forbid exfiltrating records through side channels such as deliberately
inflated error messages; the probe encoded no respondent value, prediction,
item name or score (Appendix~\ref{app:probe}).

%% ---------------------------------------------------------------------------
\section{Results}
\label{sec:results}

\subsection{Test-phase leaderboard results}

\begin{table}[H]
  \centering
  \small
  \caption{Official SimulacraBench results for team \teamid, as shown on the
  leaderboards (precision $0.01$). Development: our leaderboard entry
  (submission 1050523); Test: submission 1059398. Skill is defined so that the
  organisers' \texttt{baseline/marginal\_counts} scores $0$.}
  \label{tab:results}
  \begin{tabular}{lrrrr}
    \toprule
    Split & \texttt{skill} & \texttt{skill\_unicef} & \texttt{skill\_world\_bank} & \texttt{skill\_unhcr} \\
    \midrule
    Development & 0.48 & 0.35 & 0.23 & 0.87 \\
    Test        & 0.48 & 0.34 & 0.23 & 0.87 \\
    Organizer baseline (test) & 0.00 & 0.00 & 0.00 & 0.00 \\
    \bottomrule
  \end{tabular}
\end{table}

\subsection{Breakdown by population, survey, or question type}
The leaderboard reports one score per survey. The method does worst on World
Bank ($0.23$) and UNICEF ($0.34$), where the GIVEN block says little about most
answers, and best on UNHCR ($0.87$), where many items are gated and the skip
structure is predictable from GIVEN and parent answers. Development versions
with and without TabICLv2 fell within one rounding step of each other
(Appendix~\ref{app:results}), so we cannot say which component helps on the
real data.

%% ---------------------------------------------------------------------------
\section{Reproducibility}
\label{sec:repro}

\begin{itemize}[leftmargin=*]
  \item Code repository: \url{\CodeURL} (licence: MIT).
  \item Commit hash used for the test-phase submission: \texttt{f6b2e2002efe}
    (\texttt{main.py}, \texttt{requirements.txt}, \texttt{models.txt}; the ZIP
    has SHA-256 \texttt{\seqsplit{c0e043749d6803bdab4a0a2a0fc09dcc29898ad9ee6984be5cb9007cc07b17a3}}).
  \item Command to reproduce the test-phase submission file:
    \texttt{zip sub.zip main.py requirements.txt models.txt} at that commit.
    The predictions are computed by the grader from the released data inside
    \texttt{predict()}, so only the organisers can rerun the ZIP on the real
    data; others can run it end to end with the organisers'
    \texttt{make\_sandbox.py} and \texttt{score.py} (repository README).
  \item Environment: the grader's hosted Python 3.13 image (numpy, pandas,
    scikit-learn, torch preinstalled) plus \texttt{tabicl==2.2.0} and
    \texttt{einops==0.8.2}. All randomness derives from seed \texttt{20261112}.
    Reruns match exactly only if no wall-clock guard fires, and up to GPU
    floating-point non-determinism. Without a GPU the code skips TabICL unless
    \texttt{SB\_TABICL\_CPU=1} is set.
  \item Pretrained weights: \TabICLUsed, public on Hugging Face and still
    obtainable. \texttt{models.txt} names the repository, not a revision; its
    head has been \texttt{4dcd344e} since 2026-02-18.
\end{itemize}

%% ---------------------------------------------------------------------------
\section{Limitations}
\label{sec:limitations}

The backoff and the trees treat every GIVEN column as categorical and ignore the
order of ordinal answers, and the gate recursion assumes the declared parent is
the only reason an item is skipped. The out-of-fold estimate used to choose
weights and temperatures is optimistic, since they are chosen on the same folds,
and Development scores could not separate our versions, so per-component gains
on the real data are unknown.

\paragraph{Use of AI tools.} The code and this report were written with an AI
coding assistant (Claude) under the author's direction; every number in the
report comes from our own runs or from the cited sources.

%% ---------------------------------------------------------------------------
\bibliographystyle{plainnat}
\bibliography{refs}

%% ---------------------------------------------------------------------------
\clearpage
\appendix

\section{Full configuration}
\label{app:config}

\begin{table}[h]
  \centering
  \small
  \begin{tabular}{@{}lp{0.76\linewidth}@{}}
    \toprule
    Setting & Value \\
    \midrule
    Folds; seed & 3 per instrument; \texttt{20261112} (folds, network, trees and TabICL derive from it) \\
    Grader floor & $\varepsilon = 10^{-3}$ (used in every out-of-fold comparison) \\
    Backoff & $\alpha \in \{0.5, 2, 8, 32\}$, depth $\in \{0,\dots,4\}$, chosen out of fold; level 0 has half a count per option \\
    Gate recursion & $g(v)$ smoothed by half a pseudo-count toward the declared rule \\
    Network & one hidden layer of 256, dropout 0.25, AdamW (lr $3\cdot10^{-3}$, weight decay $10^{-4}$), batch 256, at most 60 epochs, patience 6, 3 seeds per fold \\
    Trees & learning rate 0.08, at most 200 iterations, 15 leaves, early stopping on 15\% \\
    TabICLv2 & 4 estimators, at most 20{,}000 context rows, no fine-tuning \\
    Mixing & greedy, in the order network, trees, TabICL; $w \in \{0.1, \dots, 1.0\}$ \\
    Temperature & $t \in \{0.7, 0.8, 0.9, 1.1, 1.25, 1.5\}$, applied as $p^t$ renormalised \\
    Guards (Test) & trees 2016\,s, TabICL 2160\,s, soft budget 2520\,s (Development: 560, 600, 700\,s) \\
    \bottomrule
  \end{tabular}
\end{table}

\noindent Features: all GIVEN columns of the instrument, categorical; the backoff
ranks them by mutual information with the item's answer on the training fold
and forces a GIVEN gate parent first. No prompts are used.

\section{(Optional) Additional results}
\label{app:results}

\paragraph{Development versions.} Development skill as returned for each of our
submissions. Development scores are noised and rounded to $0.01$ on the uniform
scale before conversion, so one step is about $0.0135$, $0.0150$ and $0.0193$
skill on UNICEF, World Bank and UNHCR.

\begin{center}\small
\begin{tabular}{@{}llrrrr@{}}
\toprule
Version & What changes & UNICEF & World Bank & UNHCR & Overall\\
\midrule
v4 & backoff, recursion, network, trees & 0.3499 & 0.2335 & 0.8457 & 0.4764\\
v5 & + TabICLv2 & 0.3499 & 0.2335 & 0.8650 & 0.4828\\
v5 & same ZIP, next day & 0.3499 & 0.2335 & 0.8457 & 0.4764\\
v6, v7 & TabICL time shared or given to UNHCR & 0.3499 & 0.2335 & 0.8457 & 0.4764\\
v8 & v5 + refit on all visible rows for hidden rows & 0.3499 & 0.2335 & 0.8457 & 0.4764\\
\bottomrule
\end{tabular}
\end{center}

\noindent The same v5 ZIP scored $0.8650$ and $0.8457$ on UNHCR on consecutive
days, so Development scores could not rank v4--v8; we chose the final version
(v8 with the 3600\,s budget) by the in-container out-of-fold selection itself.
v4 (no pretrained model) was the fallback ZIP and was not needed. Our own
submission record gives the unrounded Test scores of submission 1059398 as
$0.4832$ overall, $0.3419$ UNICEF, $0.2343$ World Bank and $0.8734$ UNHCR.

\paragraph{The probe.}
\label{app:probe}
The grader returns no standard output or traceback; for an exception inside
\texttt{predict()} it returns only the exception's type, and failed submissions
do not count. Since a silent fallback (no GPU, checkpoint not found, out of
time) would leave the score close to the version without TabICL, we submitted
once, in the Test phase, a probe identical to the final code that records how
far TabICL got and, after the last instrument, raises an exception whose type
encodes it: for example \texttt{ProcessLookupError} for ``checkpoint not
found'', \texttt{RecursionError} for ``ran but chosen for no item'', and
\texttt{NotImplementedError} for ``mixed into at least one item, run finished
within 75\% of the budget''. The competition terms forbid exfiltrating records
through side channels such as deliberately inflated error messages. The probe
carried no record: the exception type took one of eleven fixed values set only
by the state of our own code (GPU visible, checkpoint found and imported, how
far the TabICL fits got, whether out-of-fold selection mixed TabICL into any
item, and whether the run used less than 75\% of the budget). No respondent
value, prediction, item name or score was encoded; the exception message, which
the grader does not return, held only counts of fitted and mixed items and the
elapsed time. The answer key never enters the container (README, ``What happens
when you submit'', step~3). The probe is in the repository (\texttt{probe/}).
\textbf{Result:} \texttt{\ProbeType}, i.e.\ \ProbeMeaning.

\end{document}
